\documentclass[12pt, a4paper]{article}

% ── Packages ──────────────────────────────────────────────────────────────────
\usepackage[margin=2.5cm]{geometry}
\usepackage{amsmath, amssymb, amsthm}
\usepackage{booktabs}
\usepackage{array}
\usepackage{longtable}
\usepackage{hyperref}
\usepackage{xcolor}
\usepackage{titlesec}
\usepackage{fancyhdr}
\usepackage{graphicx}
\usepackage{float}
\usepackage{enumitem}
\usepackage{mathtools}
\usepackage{bbm}
\usepackage{caption}
\usepackage{subcaption}
\usepackage[round]{natbib}

% ── Colours ───────────────────────────────────────────────────────────────────
\definecolor{darkblue}{RGB}{0,51,102}
\definecolor{midblue}{RGB}{0,102,179}
\definecolor{lightgray}{RGB}{245,245,245}

% ── Hyperref setup ────────────────────────────────────────────────────────────
\hypersetup{
    colorlinks=true,
    linkcolor=darkblue,
    urlcolor=midblue,
    citecolor=darkblue,
    pdfauthor={Risk Management},
    pdftitle={Loan-to-Value Methodology for Margin Trading}
}

% ── Header / Footer ───────────────────────────────────────────────────────────
\pagestyle{fancy}
\fancyhf{}
\rhead{\textcolor{darkblue}{\small LTV Methodology v1.0}}
\lhead{\textcolor{darkblue}{\small Confidential}}
\cfoot{\thepage}

% ── Section formatting ────────────────────────────────────────────────────────
\titleformat{\section}{\large\bfseries\color{darkblue}}{{\thesection}}{1em}{}[\titlerule]
\titleformat{\subsection}{\normalsize\bfseries\color{midblue}}{\thesubsection}{1em}{}

% ── Theorem environments ──────────────────────────────────────────────────────
\newtheorem{definition}{Definition}[section]
\newtheorem{proposition}{Proposition}[section]
\newtheorem{remark}{Remark}[section]

% ── Custom commands ───────────────────────────────────────────────────────────
\newcommand{\LTV}{\mathrm{LTV}}
\newcommand{\E}{\mathbb{E}}
\newcommand{\Prob}{\mathbb{P}}
\newcommand{\R}{\mathbb{R}}
\newcommand{\1}{\mathbbm{1}}
\newcommand{\VaR}{\mathrm{VaR}}
\newcommand{\CVaR}{\mathrm{CVaR}}
\DeclareMathOperator*{\argmin}{arg\,min}

% ══════════════════════════════════════════════════════════════════════════════
\begin{document}

% ── Title Page ────────────────────────────────────────────────────────────────
\begin{titlepage}
    \centering
    \vspace*{2cm}
    {\color{darkblue}\rule{\linewidth}{1.5pt}}\\[0.4cm]
    {\Huge\bfseries\color{darkblue} Loan-to-Value Methodology\\[0.3cm]
    for Margin Trading}\\[0.4cm]
    {\color{darkblue}\rule{\linewidth}{1.5pt}}\\[2cm]

    {\large\textbf{Version:} 1.0}\\[0.3cm]
    {\large\textbf{Classification:} Confidential}\\[0.3cm]
    {\large\textbf{Owner:} Risk Management / Quantitative Risk}\\[0.3cm]
    {\large\textbf{Date:} \today}\\[3cm]

    \begin{abstract}
    \noindent
    This document defines the quantitative methodology governing the calculation,
    monitoring, and enforcement of Loan-to-Value (LTV) ratios within the margin
    trading product. It serves as the contractual and operational reference between
    the Platform and the external Lender, covering collateral eligibility, haircut
    derivation, LTV computation, threshold calibration, dynamic regime adjustment,
    margin call procedures, forced liquidation logic, stress testing, and governance.
    The framework is designed to be mathematically rigorous, operationally
    transparent, and aligned with applicable regulatory standards (MiFID~II,
    CRR~II, ESMA leverage guidelines).
    \end{abstract}

    \vfill
    {\small\color{gray} This document is confidential and intended solely for use
    by authorised personnel of the Platform and the Lender. Reproduction or
    distribution without written consent is prohibited.}
\end{titlepage}

% ── Table of Contents ─────────────────────────────────────────────────────────
\tableofcontents
\newpage

% ══════════════════════════════════════════════════════════════════════════════
\section{Purpose and Scope}
\label{sec:scope}

\subsection{Purpose}
This methodology document establishes the quantitative framework governing
Loan-to-Value (LTV) ratios for the margin trading product offered by the Platform.
It provides a single authoritative reference for:
\begin{itemize}[leftmargin=1.5em]
    \item the mathematical definition and computation of LTV;
    \item the derivation of collateral haircuts using risk-based methods;
    \item dynamic threshold adjustment under changing market regimes;
    \item margin call and forced liquidation procedures;
    \item stress testing and scenario analysis of the collateral portfolio;
    \item governance, backtesting, and model validation requirements.
\end{itemize}

\subsection{Scope}
This methodology applies to all margin loans extended to retail and professional
clients, all collateral held in segregated custody accounts, and all lifecycle
events affecting the loan balance or collateral portfolio. It does \emph{not} cover
unsecured lending, short-selling margin requirements, or derivatives margining.

\subsection{Regulatory Context}
The framework is aligned with:
\begin{itemize}[leftmargin=1.5em]
    \item \textbf{MiFID II / MiFIR} -- suitability and leverage requirements;
    \item \textbf{CRR II / Basel III} -- Financial Collateral Comprehensive Method (FCCM) for haircut recognition;
    \item \textbf{ESMA Guidelines on Leverage} -- retail client leverage caps by asset class;
    \item \textbf{EMIR} -- where collateral overlaps with derivative positions.
\end{itemize}

% ══════════════════════════════════════════════════════════════════════════════
\section{Definitions}
\label{sec:definitions}

\begin{definition}[Loan Amount]
The loan amount at time $t$ is denoted $L_t$ and comprises outstanding principal
plus accrued interest.
\end{definition}

\begin{definition}[Gross Market Value]
The gross market value of the collateral portfolio at time $t$ is
\[
V_t^{\mathrm{gross}} = \sum_{i=1}^{N} P_{i,t}\,Q_i,
\]
where $P_{i,t}$ is the last traded price of asset $i$ and $Q_i$ is the quantity held.
\end{definition}

\begin{definition}[Adjusted Collateral Value]
The adjusted collateral value $V_t^{\mathrm{adj}}$ is the gross market value
after application of per-asset haircuts, concentration caps, and a
portfolio-level correlation adjustment (defined in Section~\ref{sec:collateral}).
\end{definition}

\begin{definition}[Loan-to-Value Ratio]
\[
\LTV_t \;=\; \frac{L_t}{V_t^{\mathrm{adj}}}
\]
\end{definition}

\begin{definition}[Haircut]
The haircut $h_i \in [0,1)$ for asset $i$ is the proportional reduction applied to
its market value to derive its eligible collateral contribution, reflecting
liquidation risk.
\end{definition}

\begin{definition}[Concentration Adjustment]
A portfolio-level multiplicative factor $\kappa_t \in [0,1)$ applied to
$V_t^{\mathrm{adj}}$ to penalise undiversified or highly correlated collateral
portfolios.
\end{definition}

% ══════════════════════════════════════════════════════════════════════════════
\section{Collateral Eligibility and Haircut Framework}
\label{sec:collateral}

\subsection{Eligibility Criteria}
An asset qualifies as eligible collateral if and only if it satisfies \emph{all}
of the following criteria simultaneously:
\begin{enumerate}[leftmargin=1.5em]
    \item \textbf{Listing:} Traded on a recognised exchange (EU regulated market
    or equivalent third-country venue).
    \item \textbf{Liquidity:} 30-day average daily traded value (ADTV)
    $\geq \EUR{500{,}000}$.
    \item \textbf{Price Availability:} End-of-day and intraday prices available
    from a primary source and at least one independent fallback.
    \item \textbf{Instrument Type:} Equities, broad-market ETFs, government bonds,
    investment-grade corporate bonds. Explicitly excluded: structured products,
    leveraged/inverse ETFs, unlisted securities, crypto-assets.
    \item \textbf{Jurisdiction:} Issuer domiciled or listed in EEA, US, UK,
    Switzerland, Japan, Australia, or Canada.
\end{enumerate}

\subsection{Risk-Based Haircut Derivation}
\label{subsec:haircut}

Haircuts are designed to absorb the potential mark-to-market loss over the
liquidation horizon with high probability. The base haircut for asset $i$ is
derived as the Value-at-Risk of a long position over a holding period $\Delta t$:

\begin{equation}
    h_i \;=\; 1 - \frac{1}{1 + r_f\,\Delta t}\,
    \exp\!\left(-\,z_\alpha\,\hat{\sigma}_i\,\sqrt{\Delta t}\right),
    \label{eq:haircut}
\end{equation}

where:
\begin{itemize}[leftmargin=1.5em]
    \item $\hat{\sigma}_i$ -- annualised return volatility of asset $i$, estimated
    via an EWMA model with decay factor $\lambda = 0.94$ over a 252-day window:
    \[
    \hat{\sigma}_i^2 = (1-\lambda)\sum_{k=0}^{T-1}\lambda^k r_{i,t-k}^2,
    \]
    with daily log-returns $r_{i,t} = \ln(P_{i,t}/P_{i,t-1})$;
    \item $\Delta t = 5/252$ -- assumed liquidation horizon of five trading days;
    \item $z_\alpha = \Phi^{-1}(\alpha)$ -- standard normal quantile at confidence
    level $\alpha = 99\%$, i.e.\ $z_{0.99} = 2.326$;
    \item $r_f$ -- ECB deposit facility rate (annualised).
\end{itemize}

\begin{remark}
    For assets with fat-tailed return distributions, the Normal quantile may
    underestimate tail risk. A Cornish--Fisher expansion can be applied:
    \[
    z_\alpha^{CF} = z_\alpha + \frac{z_\alpha^2 - 1}{6}\,\hat{\gamma}_1
    + \frac{z_\alpha^3 - 3z_\alpha}{24}\,\hat{\gamma}_2
    - \frac{2z_\alpha^3 - 5z_\alpha}{36}\,\hat{\gamma}_1^2,
    \]
    where $\hat{\gamma}_1$ and $\hat{\gamma}_2$ are the estimated skewness and
    excess kurtosis of the return series.
\end{remark}

\subsubsection*{Liquidity-Adjusted Haircut}

An additional bid-ask spread adjustment is applied for assets with elevated
transaction costs:
\begin{equation}
    h_i^{\mathrm{liq}} = h_i + \max\!\left(0,\;\frac{s_i - s^*}{2}\right),
    \label{eq:liq_haircut}
\end{equation}
where $s_i$ is the observed time-weighted average bid-ask spread and $s^* =
0.10\%$ is the reference spread for liquid large-cap equities.

\subsubsection*{Indicative Haircut Table}

Table~\ref{tab:haircuts} presents the operationally applied haircut schedule,
derived from equation~\eqref{eq:haircut} and reviewed quarterly.

\begin{table}[H]
\centering
\caption{Indicative Haircut Schedule by Asset Class}
\label{tab:haircuts}
\begin{tabular}{llr}
\toprule
\textbf{Asset Class} & \textbf{Subcategory} & \textbf{Base Haircut $h_i$} \\
\midrule
Government Bonds & AAA--AA, maturity $< 1$Y     & 0.5\% \\
                 & AAA--AA, 1--5Y                & 2.0\% \\
                 & AAA--AA, 5--10Y               & 4.0\% \\
                 & A--BBB, any maturity          & 6--12\% \\
IG Corporate Bonds & 0--5Y                       & 8.0\% \\
                   & 5Y+                         & 12.0\% \\
Large-Cap Equities & MSCI World, ADTV $> \EUR{5M}$ & 20.0\% \\
Mid-Cap Equities   & ADTV $\EUR{500K}$--$\EUR{5M}$ & 30.0\% \\
Small-Cap Equities & ADTV $< \EUR{500K}$ (min)   & 40.0\% \\
ETFs (broad)       & Major index replication     & 20.0\% \\
ETFs (thematic)    & Sector/factor exposure      & 30.0\% \\
\bottomrule
\end{tabular}
\end{table}

\subsection{Concentration Limits and Portfolio Adjustment}

\subsubsection*{Single-Name Cap}
The eligible collateral value attributed to any single issuer is capped at 20\%
of the unadjusted collateral value. Let $w_i = P_{i,t}Q_i(1-h_i) /
\sum_j P_{j,t}Q_j(1-h_j)$ be the weight of issuer $i$. The concentration-adjusted
contribution is:
\[
C_i = P_{i,t}\,Q_i\,(1 - h_i) \cdot \min\!\left(1,\;\frac{0.20}{w_i}\right).
\]

\subsubsection*{Portfolio Correlation Adjustment}
Let $\bar{\rho}_t$ denote the average pairwise Pearson correlation of daily
log-returns across all collateral positions, estimated over a 60-day rolling window.
The correlation adjustment is:
\begin{equation}
    \kappa(\bar{\rho}_t) =
    \begin{cases}
        0 & \bar{\rho}_t \leq 0.4 \\[4pt]
        0.05\cdot\dfrac{\bar{\rho}_t - 0.4}{0.6} & 0.4 < \bar{\rho}_t \leq 1
    \end{cases}
    \label{eq:kappa}
\end{equation}
This adjustment linearly ramps from 0\% to 5\% as average pairwise correlation
rises from 0.4 to 1.0, reflecting the reduced diversification benefit in stress.

% ══════════════════════════════════════════════════════════════════════════════
\section{LTV Calculation Methodology}
\label{sec:ltv_calc}

\subsection{Core LTV Formula}

The adjusted collateral value incorporating all adjustments is:
\begin{equation}
    V_t^{\mathrm{adj}} = \bigl(1 - \kappa(\bar{\rho}_t)\bigr)
    \sum_{i=1}^{N} C_i(t),
    \label{eq:Vadj}
\end{equation}
and the LTV is:
\begin{equation}
    \LTV_t = \frac{L_t}{V_t^{\mathrm{adj}}}.
    \label{eq:ltv}
\end{equation}

\subsection{Loan Balance Dynamics}

The loan balance accrues continuously at the margin lending rate $r_m$
(floating: EURIBOR 3M + contractual spread), reduced by any client repayments:
\begin{equation}
    L_t = L_0\,e^{r_m t} - \sum_{k:\,t_k \leq t} R_k\,e^{r_m(t - t_k)},
    \label{eq:loan}
\end{equation}
where $R_k$ is the repayment made at time $t_k$.

\subsection{Calculation Frequency and Price Sources}

\begin{table}[H]
\centering
\caption{LTV Calculation Schedule}
\begin{tabular}{lll}
\toprule
\textbf{Calculation} & \textbf{Frequency} & \textbf{Trigger} \\
\midrule
Real-time LTV    & Continuous (tick)          & Any price update \\
Intraday Snap    & Every 15 min (market hrs)  & Scheduled \\
End-of-Day LTV   & Once, post close           & Official close prices \\
Stress LTV       & Daily, post close          & Scheduled \\
Regime LTV       & Daily, pre-open            & Regime classification \\
\bottomrule
\end{tabular}
\end{table}

\noindent\textbf{Price source hierarchy:}
\begin{enumerate}[leftmargin=1.5em]
    \item Primary: Exchange last-traded price (real-time feed).
    \item Fallback 1: Bloomberg/Refinitiv composite.
    \item Fallback 2: Prior closing price with an additional 5\% liquidity haircut.
    \item If no valid price for $>2$ consecutive trading days: zero collateral value.
\end{enumerate}

% ══════════════════════════════════════════════════════════════════════════════
\section{Dynamic Threshold Adjustment}
\label{sec:dynamic}

\subsection{Motivation}

Static LTV thresholds calibrated under average market conditions may be
insufficient during periods of elevated volatility or structural regime change.
A dynamic threshold framework adjusts the warning and liquidation LTVs in
real time, ensuring that the probability of the loan becoming uncovered remains
below the target tolerance $p^*$ under all market regimes.

\subsection{Volatility Regime Classification}
\label{subsec:regime}

Market regime is classified daily using a two-state Hidden Markov Model (HMM)
estimated on the VSTOXX index (or CBOE VIX for US-dominated portfolios).
Let the latent state $S_t \in \{0, 1\}$ denote \emph{normal} and \emph{stress}
regimes respectively. The observation equation is:
\[
x_t = \mu_{S_t} + \sigma_{S_t}\,\varepsilon_t, \qquad \varepsilon_t \sim \mathcal{N}(0,1),
\]
where $x_t = \ln(\text{VSTOXX}_t)$ and the parameters
$\Theta = \{\mu_0, \sigma_0, \mu_1, \sigma_1, p_{00}, p_{11}\}$ are estimated
via the Baum--Welch (EM) algorithm.

The smoothed state probability is:
\[
\pi_t^{(1)} = \Prob(S_t = 1 \mid x_{1:t};\,\hat{\Theta}),
\]
computed via the forward--backward algorithm.

In practice, a simpler but robust \textbf{three-regime classification} based on
realised volatility percentiles is operationally preferred:

\begin{table}[H]
\centering
\caption{Volatility Regime Classification}
\label{tab:regime}
\begin{tabular}{llll}
\toprule
\textbf{Regime} & \textbf{Condition} & \textbf{VSTOXX Level} & \textbf{Label} \\
\midrule
Normal  & $\sigma_{p,5d} \leq \bar{\sigma}_p + 1\hat{s}$ & $<$ 20 & \textsc{Normal} \\
Elevated & $\bar{\sigma}_p + 1\hat{s} < \sigma_{p,5d} \leq \bar{\sigma}_p + 2\hat{s}$ & 20--30 & \textsc{Elevated} \\
Stress  & $\sigma_{p,5d} > \bar{\sigma}_p + 2\hat{s}$ & $>$ 30 & \textsc{Stress} \\
\bottomrule
\end{tabular}
\end{table}

\noindent where $\sigma_{p,5d}$ is the 5-day realised portfolio volatility,
$\bar{\sigma}_p$ is the long-run mean, and $\hat{s}$ is the long-run standard
deviation of 5-day realised volatility, both estimated over a 2-year rolling window.

\subsection{Dynamic Threshold Formula}

The regime-adjusted liquidation threshold is defined such that:
\begin{equation}
    \Prob\!\left(\Delta V_t > V_t^{\mathrm{adj}}\bigl(1 - \LTV^{\mathrm{liq}}_t\bigr)\right) \leq p^*,
    \label{eq:threshold_condition}
\end{equation}
which, under a Normal approximation with portfolio volatility $\hat{\sigma}_{p,t}$
and liquidation horizon $\Delta t$, yields:

\begin{equation}
    \LTV^{\mathrm{liq}}_t = 1 - z_{1-p^*}\,\hat{\sigma}_{p,t}\,\sqrt{\Delta t},
    \label{eq:dynamic_liq}
\end{equation}

where $\hat{\sigma}_{p,t}$ is updated daily. To avoid excessive threshold
oscillation, the dynamic threshold is smoothed via an exponential moving average
with half-life $\tau = 5$ trading days:

\begin{equation}
    \widetilde{\LTV}^{\mathrm{liq}}_t
    = \beta\,\LTV^{\mathrm{liq}}_t + (1-\beta)\,\widetilde{\LTV}^{\mathrm{liq}}_{t-1},
    \qquad \beta = 1 - e^{-\ln 2 / \tau}.
    \label{eq:threshold_smooth}
\end{equation}

Hard floors and ceilings are applied to ensure operational bounds:
\begin{equation}
    \LTV^{\mathrm{liq},\star}_t
    = \min\!\left(\LTV^{\mathrm{liq},\max},\;
    \max\!\left(\LTV^{\mathrm{liq},\min},\;
    \widetilde{\LTV}^{\mathrm{liq}}_t\right)\right).
    \label{eq:threshold_bounds}
\end{equation}

\subsection{Regime-Dependent Threshold Schedule}

Table~\ref{tab:thresholds} summarises the operationally applied thresholds across
regimes, derived from equation~\eqref{eq:dynamic_liq} for a representative
diversified equity portfolio ($p^* = 0.1\%$):

\begin{table}[H]
\centering
\caption{Regime-Dependent LTV Threshold Schedule}
\label{tab:thresholds}
\begin{tabular}{lcccc}
\toprule
\textbf{Regime} & $\hat{\sigma}_{p,t}$ (ann.) & \textbf{Initial LTV} & \textbf{Warning LTV} & \textbf{Liquidation LTV} \\
\midrule
Normal   & $\approx 15\%$ & 60\% & 70\% & 80\% \\
Elevated & $\approx 25\%$ & 55\% & 65\% & 75\% \\
Stress   & $\approx 40\%$ & 45\% & 55\% & 65\% \\
\bottomrule
\end{tabular}
\end{table}

\noindent In the Stress regime, newly originated loans must satisfy the lower
Initial LTV. Existing loans are not force-cured solely due to a regime transition,
but receive an accelerated margin call cure deadline.

\subsection{Asymmetric Regime Transitions}

To avoid threshold whipsawing around regime boundaries, asymmetric persistence
rules are applied:
\begin{itemize}[leftmargin=1.5em]
    \item \textbf{Upgrade} (Stress $\to$ Elevated or Normal): requires 5 consecutive
    trading days in the lower-volatility regime before thresholds are relaxed.
    \item \textbf{Downgrade} (Normal/Elevated $\to$ Stress): immediate threshold
    tightening on the day of classification.
\end{itemize}

% ══════════════════════════════════════════════════════════════════════════════
\section{Margin Call Process}
\label{sec:margin_call}

\subsection{Trigger and Notification}

A margin call is triggered when:
\[
\LTV_t \geq \LTV^{\mathrm{warn},\star}_t.
\]
The notification is issued within 15 minutes of trigger detection and specifies
the current LTV, the warning threshold, the required cure amount $\Delta$, and the
cure deadline.

\subsection{Cure Amount Calculation}

The cure target is restoration to the regime-adjusted initial LTV,
$\LTV^{\mathrm{target}} = \LTV^{\mathrm{init},\star}_t$. Three cure mechanisms
are available:

\subsubsection*{Option A -- Cash Deposit / Partial Repayment}
\[
\Delta^{\mathrm{cash}} = L_t - \LTV^{\mathrm{target}}\cdot V_t^{\mathrm{adj}}.
\]

\subsubsection*{Option B -- Additional Collateral}
\[
\Delta^{\mathrm{collateral}}
= \frac{L_t - \LTV^{\mathrm{target}}\cdot V_t^{\mathrm{adj}}}
       {\LTV^{\mathrm{target}}\cdot(1 - h_{\mathrm{new}})},
\]
where $h_{\mathrm{new}}$ is the haircut of the newly deposited collateral.

\subsubsection*{Option C -- Partial Position Reduction}
\[
\Delta^{\mathrm{reduce}}
= \frac{L_t - \LTV^{\mathrm{target}}\cdot V_t^{\mathrm{adj}}}
       {1 - \LTV^{\mathrm{target}}\cdot(1 - h_{\mathrm{sold}})},
\]
where $h_{\mathrm{sold}}$ is the haircut of the sold collateral.

\subsection{Cure Period}
\begin{table}[H]
\centering
\begin{tabular}{lll}
\toprule
\textbf{Client Type} & \textbf{Normal Regime} & \textbf{Stress Regime} \\
\midrule
Retail       & T+1, 12:00 CET & Same-day, 4 hours \\
Professional & T+0, 4 hours   & T+0, 2 hours \\
\bottomrule
\end{tabular}
\end{table}

% ══════════════════════════════════════════════════════════════════════════════
\section{Forced Liquidation Process}
\label{sec:liquidation}

\subsection{Trigger Conditions}
Forced liquidation is initiated when \emph{any} of the following conditions hold:
\begin{enumerate}[leftmargin=1.5em]
    \item $\LTV_t \geq \LTV^{\mathrm{liq},\star}_t$ at any intraday snapshot;
    \item Cure period expires without sufficient cure action;
    \item $\LTV_t^{\mathrm{stress}} \geq \LTV^{\mathrm{liq},\star}_t$ on two
    consecutive business days (pre-emptive trigger).
\end{enumerate}

\subsection{Liquidation Sequencing}

Positions are ranked by liquidity score:
\[
\ell_i = \frac{\mathrm{ADTV}_i}{P_{i,t}\,Q_i},
\]
and liquidated in descending order of $\ell_i$. The quantity of asset $i$ to sell is:
\[
Q_i^{\mathrm{sell}}
= \min\!\left(Q_i,\;
  \frac{\Delta^{\mathrm{reduce}}}{P_{i,t}(1 - s_i/2)}\right),
\]
where $s_i/2$ is the estimated half-spread. The algorithm iterates until
$\LTV_t \leq \LTV^{\mathrm{target}}$.

\subsection{Market Impact Estimation}

For positions where $Q_i^{\mathrm{sell}} / \mathrm{ADTV}_i > 10\%$, market impact
is estimated using the Almgren--Chriss permanent impact model:
\begin{equation}
    MI_i = \eta\,\sigma_i\,\frac{Q_i^{\mathrm{sell}}}{\mathrm{ADTV}_i}\,P_{i,t},
    \label{eq:market_impact}
\end{equation}
where $\eta \in [0.1, 0.3]$ is the market impact coefficient, estimated empirically.
The total liquidation shortfall estimate is:
\[
\mathrm{TIS} = \sum_{i} MI_i + \sum_{i} P_{i,t}\,Q_i^{\mathrm{sell}}\,\frac{s_i}{2}.
\]
This shortfall is factored into the target cure amount.

\subsection{Proceeds Waterfall}

Liquidation proceeds are applied in strict priority:
\begin{enumerate}[leftmargin=1.5em]
    \item Outstanding loan principal $L_t$;
    \item Accrued interest;
    \item Platform liquidation fee;
    \item Residual returned to client.
\end{enumerate}

% ══════════════════════════════════════════════════════════════════════════════
\section{Stress Testing Framework}
\label{sec:stress}

\subsection{Overview}

The stress testing framework serves three distinct purposes:
\begin{enumerate}[leftmargin=1.5em]
    \item \textbf{Pre-emptive monitoring:} Daily computation of Stress LTV to
    detect latent vulnerabilities before mark-to-market LTV breaches thresholds.
    \item \textbf{Threshold calibration:} Stress scenarios inform the setting of
    regime-adjusted thresholds (Section~\ref{sec:dynamic}).
    \item \textbf{Lender capital planning:} Portfolio-level stress results are
    reported to the Lender for concentration risk and capital adequacy assessment.
\end{enumerate}

\subsection{Single-Asset Stress LTV}

For each client portfolio, a Stress LTV is computed by applying instantaneous
shocks $\delta_i$ to all collateral positions:

\begin{equation}
    V_t^{\mathrm{stress}} = \sum_{i=1}^{N} P_{i,t}\,Q_i\,(1-h_i)\,(1-\delta_i),
    \qquad
    \LTV_t^{\mathrm{stress}} = \frac{L_t}{V_t^{\mathrm{stress}}}.
    \label{eq:stress_ltv}
\end{equation}

The stress shock $\delta_i$ is defined as:
\[
\delta_i = \max\!\left(\delta_i^{\mathrm{hist}},\;\delta_i^{\mathrm{floor}}\right),
\]
where:
\begin{itemize}[leftmargin=1.5em]
    \item $\delta_i^{\mathrm{hist}}$ is the 99th percentile of the empirical
    distribution of 5-day absolute log-returns of asset $i$ over a 5-year lookback:
    \[
    \delta_i^{\mathrm{hist}} = \widehat{F}_i^{-1}(0.99),\qquad
    \widehat{F}_i = \text{empirical CDF of }|r_{i,t:t+5\Delta}|;
    \]
    \item $\delta_i^{\mathrm{floor}} = 15\%$ for equities, $5\%$ for investment-grade
    bonds (regulatory floor).
\end{itemize}

\subsection{Historical Scenario Analysis}

The collateral portfolio is revalued under the following named historical
stress scenarios, replaying actual observed returns:

\begin{table}[H]
\centering
\caption{Named Historical Stress Scenarios}
\label{tab:scenarios}
\begin{tabular}{llr}
\toprule
\textbf{Scenario} & \textbf{Period} & \textbf{Approx.\ Peak Equity Drawdown} \\
\midrule
Global Financial Crisis   & Oct 2008 -- Mar 2009 & $-50\%$ \\
European Sovereign Crisis & Aug 2011 -- Sep 2011 & $-22\%$ \\
COVID-19 Shock            & Feb 2020 -- Mar 2020 & $-35\%$ \\
Rate Shock 2022           & Jan 2022 -- Oct 2022 & $-25\%$ (equities), $-20\%$ (bonds) \\
Flash Crash               & Aug 24, 2015         & $-11\%$ intraday \\
\bottomrule
\end{tabular}
\end{table}

For each scenario $s$ with path $\{r_{i,t}^{(s)}\}$, the portfolio trajectory is:
\[
V_\tau^{(s)} = \sum_{i=1}^{N} P_{i,0}\,Q_i\,(1-h_i)
\prod_{t=1}^{\tau}\bigl(1 + r_{i,t}^{(s)}\bigr),
\qquad \LTV_\tau^{(s)} = \frac{L_\tau}{V_\tau^{(s)}},
\]
and the scenario breach indicator is:
\[
\mathbbm{1}^{(s)}_{\mathrm{breach}}
= \max_{\tau}\,\mathbbm{1}\!\left[\LTV_\tau^{(s)} \geq \LTV^{\mathrm{liq},\star}\right].
\]

\subsection{Hypothetical (Adverse) Scenario Analysis}

In addition to historical replays, the following hypothetical shocks are applied:

\begin{enumerate}[leftmargin=1.5em]
    \item \textbf{Parallel equity drawdown:} All equity positions simultaneously
    lose $X\%$ over one trading day. $X$ is swept from 10\% to 50\%.
    \item \textbf{Rates shock:} Instantaneous parallel shift of $+200$bps applied
    to bond positions (duration-adjusted price impact).
    \item \textbf{Correlation breakdown:} All pairwise correlations set to 1.0;
    haircuts are recomputed accordingly and $\kappa = 5\%$ applied.
    \item \textbf{Liquidity seizure:} All ADTV values reduced by 80\%; haircuts
    recalculated via equation~\eqref{eq:liq_haircut} with inflated spreads.
    \item \textbf{FX shock:} Non-EUR assets subject to an additional 15\% adverse
    currency move.
\end{enumerate}

\subsection{Reverse Stress Testing}
\label{subsec:reverse_stress}

Reverse stress testing identifies the minimum market shock $\delta^*$ required
to breach the liquidation threshold, i.e.\ to make the loan uncovered. This is
formulated as a constrained optimisation:

\begin{equation}
    \delta^* = \argmin_{\boldsymbol{\delta}\,\geq\,0}
    \;\|\boldsymbol{\delta}\|_2
    \quad\text{s.t.}\quad
    L_t \;\geq\; V_t^{\mathrm{adj}}\bigl(\boldsymbol{\delta}\bigr),
    \label{eq:reverse_stress}
\end{equation}
where $\|\boldsymbol{\delta}\|_2 = \sqrt{\sum_i \delta_i^2}$ is the Euclidean
norm of the shock vector and:
\[
V_t^{\mathrm{adj}}(\boldsymbol{\delta})
= \bigl(1 - \kappa\bigr)
\sum_{i=1}^{N} P_{i,t}\,Q_i\,(1-h_i)\,(1-\delta_i).
\]

\noindent The closed-form solution, assuming no binding concentration constraints, is:

\begin{equation}
    \delta_i^* = \delta^*_{\mathrm{uniform}}
    = 1 - \frac{V_t^{\mathrm{adj}}(0)}{L_t / (1-\kappa)}
    \cdot \frac{1}{N},
    \label{eq:reverse_uniform}
\end{equation}
for a uniform shock. For a risk-weighted shock (proportional to volatility):
\[
\delta_i^* = \lambda^*\,\sigma_i,
\qquad
\lambda^* = \frac{V_t^{\mathrm{adj}}(0) - L_t/(1-\kappa)}
                 {\sum_j P_{j,t}\,Q_j\,(1-h_j)\,\sigma_j}.
\]
The reverse stress result $\delta^*$ is reported alongside each daily
Stress LTV as a margin of safety metric.

\subsection{Portfolio-Level VaR and CVaR of LTV}

For the purpose of Lender reporting, a distribution of $\LTV_{t+\Delta t}$ is
estimated via Monte Carlo simulation:

\begin{enumerate}[leftmargin=1.5em]
    \item Draw $M = 10{,}000$ correlated return scenarios from a multivariate
    $t$-distribution with estimated covariance $\hat{\Sigma}$ and degrees of
    freedom $\nu$ (fitted via maximum likelihood):
    \[
    \mathbf{r}_{t+\Delta t}^{(m)} \sim t_\nu(\mathbf{0},\,\hat{\Sigma}\,\Delta t),
    \qquad m = 1,\ldots,M.
    \]
    \item Revalue the portfolio under each scenario:
    \[
    V_{t+\Delta t}^{(m)}
    = \sum_{i=1}^{N} P_{i,t}\,e^{r_{i,t+\Delta t}^{(m)}}\,Q_i\,(1-h_i).
    \]
    \item Compute the simulated LTV:
    \[
    \LTV_{t+\Delta t}^{(m)} = \frac{L_{t+\Delta t}}{V_{t+\Delta t}^{(m)}}.
    \]
    \item Report:
    \begin{align*}
    \VaR_{\alpha}(\LTV) &= \hat{F}_{\LTV}^{-1}(\alpha),\\
    \CVaR_{\alpha}(\LTV) &= \frac{1}{1-\alpha}\int_\alpha^1 \hat{F}_{\LTV}^{-1}(u)\,du.
    \end{align*}
\end{enumerate}

% ══════════════════════════════════════════════════════════════════════════════
\section{Collateral Revaluation and Corporate Actions}
\label{sec:corp_actions}

\begin{table}[H]
\centering
\caption{Corporate Action Treatment}
\begin{tabular}{p{3.5cm}p{9cm}}
\toprule
\textbf{Event} & \textbf{Treatment} \\
\midrule
Cash Dividend      & Ex-date price drop reflected via feed; dividend cash credited and may offset loan \\
Stock Split        & $Q_i \to k Q_i$, $P_i \to P_i/k$; LTV unchanged if processed correctly \\
Spin-off           & New entity assessed for eligibility; if ineligible, zero collateral value from ex-date \\
M\&A (cash)        & Position removed on settlement; proceeds applied to loan \\
M\&A (stock)       & New shares reassessed; LTV recalculated \\
Delisting          & Immediate zero collateral value; margin call if $\LTV \geq \LTV^{\mathrm{warn}}$ \\
Trading Halt       & Fallback price applied; zero value if halt $> 2$ days \\
\bottomrule
\end{tabular}
\end{table}

\subsection{FX Revaluation}
For assets denominated in currency $c \neq$ EUR, a combined haircut is applied:
\[
h_i^{\mathrm{total}} = 1 - (1-h_i)(1-h_{\mathrm{FX}(c)}),
\]
where $h_{\mathrm{FX}(c)}$ is derived from equation~\eqref{eq:haircut} applied
to the $c$/EUR exchange rate with a 5-day liquidation horizon.

% ══════════════════════════════════════════════════════════════════════════════
\section{Governance and Validation}
\label{sec:governance}

\subsection{Backtesting and the Kupiec POF Test}

Haircut adequacy is validated quarterly. For each asset class $j$, the test
statistic is:
\begin{equation}
    LR_{\mathrm{POF}}^{(j)}
    = -2\ln\!\left[
        \frac{(1-p)^{T-x}\,p^x}{(1-\hat{p})^{T-x}\,\hat{p}^x}
    \right] \;\xrightarrow{d}\; \chi^2(1),
    \label{eq:kupiec}
\end{equation}
where $p$ is the target exceedance rate ($p = 1\%$ for $\alpha = 99\%$),
$T$ is the sample size, $x$ is the number of observed haircut breaches, and
$\hat{p} = x/T$. Haircuts are flagged for revision if $LR > \chi^2_{0.95}(1) = 3.84$.

Additionally, the Christoffersen conditional coverage test is applied to assess
independence of exceedances:
\[
LR_{\mathrm{CC}} = LR_{\mathrm{POF}} + LR_{\mathrm{ind}} \;\sim\; \chi^2(2).
\]

\subsection{Review Schedule}

\begin{table}[H]
\centering
\begin{tabular}{ll}
\toprule
\textbf{Activity} & \textbf{Frequency} \\
\midrule
Haircut backtesting (Kupiec + Christoffersen) & Quarterly \\
Regime model recalibration (HMM)             & Semi-annual \\
Threshold calibration review                 & Quarterly \\
Stress scenario update                       & Annual (or post-crisis) \\
Full methodology review                      & Annual \\
Lender portfolio LTV report                  & Daily \\
Lender stress report                         & Weekly \\
\bottomrule
\end{tabular}
\end{table}

% ══════════════════════════════════════════════════════════════════════════════
\appendix

\section{Notation Summary}
\label{app:notation}

\begin{longtable}{cl}
\toprule
\textbf{Symbol} & \textbf{Description} \\
\midrule
\endhead
$L_t$                        & Loan amount at time $t$ \\
$V_t^{\mathrm{gross}}$       & Gross market value of collateral at $t$ \\
$V_t^{\mathrm{adj}}$         & Adjusted collateral value at $t$ \\
$V_t^{\mathrm{stress}}$      & Stressed collateral value at $t$ \\
$P_{i,t}$                    & Price of asset $i$ at $t$ \\
$Q_i$                        & Quantity held of asset $i$ \\
$h_i$                        & Haircut for asset $i$ \\
$h_i^{\mathrm{liq}}$         & Liquidity-adjusted haircut for asset $i$ \\
$h_{\mathrm{FX}(c)}$         & FX haircut for currency $c$ \\
$\hat{\sigma}_i$             & EWMA annualised volatility of asset $i$ \\
$\hat{\sigma}_{p,t}$         & Portfolio volatility at $t$ \\
$\bar{\rho}_t$               & Average pairwise correlation at $t$ \\
$\kappa(\bar{\rho}_t)$       & Portfolio correlation adjustment \\
$S_t$                        & Volatility regime state at $t$ \\
$\pi_t^{(1)}$                & Smoothed stress-regime probability \\
$\delta_i$                   & Stress shock for asset $i$ \\
$\delta^*$                   & Reverse stress break-even shock \\
$\ell_i$                     & Liquidity score of asset $i$ \\
$r_m$                        & Margin loan rate \\
$r_f$                        & Risk-free rate \\
$\Delta t$                   & Liquidation horizon (years) \\
$z_\alpha$                   & Standard normal quantile at level $\alpha$ \\
$\lambda$                    & EWMA decay factor \\
$\eta$                       & Almgren--Chriss impact coefficient \\
$p^*$                        & Target loss probability tolerance \\
$\LTV^{\mathrm{init}}$       & Maximum initial LTV \\
$\LTV^{\mathrm{warn}}$       & Warning LTV threshold \\
$\LTV^{\mathrm{liq}}$        & Liquidation LTV threshold \\
$\LTV^{\mathrm{liq},\star}_t$& Dynamic (regime-adjusted) liquidation threshold \\
\bottomrule
\end{longtable}

\section{Worked Example: Dynamic Threshold and Stress LTV}
\label{app:example}

\subsection*{Portfolio}

\begin{itemize}[leftmargin=1.5em]
    \item 1,000 shares large-cap equity: $P = \EUR{80}$, $h = 20\%$, $\sigma = 18\%$ p.a.
    \item \EUR{20,000} nominal 3Y AA government bond: price 101\%, $h = 2\%$, $\sigma = 3\%$ p.a.
    \item 500 shares mid-cap equity: $P = \EUR{40}$, $h = 30\%$, $\sigma = 28\%$ p.a.
\end{itemize}

\[
V^{\mathrm{adj}} = 64{,}000 + 19{,}796 + 14{,}000 = \EUR{97{,}796},
\qquad L = \EUR{55{,}000},
\qquad \LTV = 56.2\%.
\]

\subsection*{Portfolio Volatility}
Assuming correlations $\rho_{12} = 0.05$, $\rho_{13} = 0.60$, $\rho_{23} = 0.05$,
portfolio weights (by adjusted value) $w_1 = 65.4\%$, $w_2 = 20.2\%$,
$w_3 = 14.3\%$:

\[
\hat{\sigma}_p^2 = \mathbf{w}^\top \Sigma \mathbf{w} \approx 0.0256
\;\Rightarrow\; \hat{\sigma}_p \approx 16.0\%\ \text{p.a.}
\]

\subsection*{Dynamic Liquidation Threshold (Normal Regime)}
\[
\LTV^{\mathrm{liq}} = 1 - 3.09 \times 0.160 \times \sqrt{5/252}
= 1 - 0.1754 = 82.5\%.
\]
After smoothing and applying bounds: $\LTV^{\mathrm{liq},\star} = 80\%$.

\subsection*{Stress LTV}
$\delta_1^{\mathrm{hist}} = 22\%$, $\delta_2^{\mathrm{hist}} = 5\%$,
$\delta_3^{\mathrm{hist}} = 28\%$:
\[
V^{\mathrm{stress}}
= 64{,}000 \times (1-0.22) + 19{,}796 \times (1-0.05) + 14{,}000 \times (1-0.28)
= 49{,}920 + 18{,}806 + 10{,}080
= \EUR{78{,}806}.
\]
\[
\LTV^{\mathrm{stress}} = \frac{55{,}000}{78{,}806} = \textbf{69.8\%}.
\]
The Stress LTV is below $\LTV^{\mathrm{liq},\star} = 80\%$: \textbf{no pre-emptive trigger}.

\subsection*{Reverse Stress Break-Even}
The uniform shock required to make $L \geq V^{\mathrm{adj}}$:
\[
\delta^* = 1 - \frac{L}{V^{\mathrm{adj}}} = 1 - \frac{55{,}000}{97{,}796} = 43.8\%.
\]
A simultaneous 43.8\% uniform drawdown across all positions would render the
loan uncovered. This exceeds the GFC peak drawdown of $\approx 35\%$, indicating
adequate stress resilience for this portfolio.

\end{document}